\section*{Capítulo \ref{ch:g12:early-earth} --- La Tierra primitiva y su atmósfera}

\begin{solution}{exo:g12:early-earth:1}
(a) la \omterm{def:g12:early-earth:primitive-atmosphere}{desgasificación}; (b) un \omterm{def:g12:early-earth:stromatolite}{estromatolito}; (c) una \omterm{def:g12:early-earth:bif}{formación de hierro
bandeado}; (d) las \omterm{def:g12:early-earth:red-beds}{capas rojas}; (e) el \omterm{def:g12:early-earth:goe}{Gran Evento de Oxidación}.
\end{solution}

\begin{solution}{exo:g12:early-earth:2}
Los circones que muestran agua líquida (4400 millones de años); los primeros
\omterm{def:g12:early-earth:stromatolite}{estromatolitos} del Pilbara (3430); las \omterm{def:g12:early-earth:bif}{formaciones de hierro bandeado} de
Hamersley (más de 2500); el final del \omterm{def:g12:early-earth:goe}{Gran Evento de Oxidación} (2060); el
final de la mayoría de las \omterm{def:g12:early-earth:bif}{formaciones de hierro bandeado} (1800).
\end{solution}

\begin{solution}{exo:g12:early-earth:3}
En $t = 0$, $L = L_\odot/(1 + 0.4) = L_\odot/1.4 \approx 0.714\,L_\odot$:
alrededor de un \qty{29}{\percent} menos brillante, casi un tercio menos que hoy.
\end{solution}

\begin{solution}{exo:g12:early-earth:4}
Que el aire contenía entonces casi nada de oxígeno. La pirita es destruida por
el oxígeno; sus granos solo podían ser arrastrados por los ríos y depositarse
intactos si el aire y el agua de la superficie carecían de él.
\end{solution}

\begin{solution}{exo:g12:early-earth:5}
El vapor de agua, que se condensó y formó el océano, y el dióxido de carbono,
que permaneció en el aire en cantidades mucho mayores que hoy y calentó la
Tierra por su efecto invernadero.
\end{solution}

\begin{solution}{exo:g12:early-earth:6}
Hace 3000 millones de años: $1 - t/t_\odot = 3/4.567 \approx 0.657$, $L \approx
1/1.263 \approx 0.792\,L_\odot$, $T \approx 255 \times 0.792^{1/4} \approx
\qty{241}{K}$. Hace 2450 millones de años: $1 - t/t_\odot \approx 0.536$, $L
\approx 1/1.215 \approx 0.823\,L_\odot$, $T \approx 255 \times 0.823^{1/4}
\approx \qty{243}{K}$.
\end{solution}

\begin{solution}{exo:g12:early-earth:7}
\qty{1}{kg} de hierro son $1000/55.85 \approx 17.9$ moles; una molécula de
oxígeno capta cuatro iones de hierro, así que $17.9/4 \approx 4.48$ moles de
\ce{O2}, es decir, $4.48 \times 32.00 \approx \qty{143}{g}$ de oxígeno.
\end{solution}

\begin{solution}{exo:g12:early-earth:8}
(1) $2.46 - 2.06 = 0.40$ miles de millones de años, 400 millones de años.
(2) De hace 4540 a hace 2460 millones de años: $2.08/4.54 \approx \qty{46}{\percent}$ de la historia de la Tierra. (3) Unos $3.43 - 2.46 \approx 1$ mil millones de años: el oxígeno que producían los microbios se
consumió durante mucho tiempo antes de poder permanecer en el aire.
\end{solution}

\begin{solution}{exo:g12:early-earth:9}
La \omterm{def:g12:early-earth:bif}{formación de hierro bandeado} muestra un mar profundo sin oxígeno, que podía
llevar hierro disuelto, bajo un mar poco profundo e iluminado con oxígeno
suficiente para precipitarlo. La capa roja, depositada en tierra firme,
muestra oxígeno en el aire y en el agua de la superficie, suficiente para
oxidar el hierro de sus granos.
\end{solution}

\begin{solution}{exo:g12:early-earth:10}
En Hamelin Pool, bahía Shark, Australia Occidental, cuya agua es dos veces más
salada que el agua de mar normal. Un tapete de microbios, sobre todo
cianobacterias, atrapa y fija granos de sedimento y hace crecer carbonato;
encima crece un tapete nuevo, y así sucesivamente, capa sobre capa.
\end{solution}

\begin{solution}{exo:g12:early-earth:11}
(1) $21 \times 10^{-5} = \qty{0.00021}{\percent}$. (2) $21/10 =
\qty{2.1}{\percent}$.
\end{solution}

\begin{solution}{exo:g12:early-earth:12}
La \omterm{def:g10:habitable-planet:equilibrium-temperature}{temperatura de equilibrio} no cambia si $L(1 - A')$ es igual a $1 - A$:
$1 - A' = 0.706/0.741 \approx 0.953$, así que $A' \approx 0.047$. La Tierra
habría tenido que reflejar menos del \qty{5}{\percent} de la luz solar, casi
como un cuerpo negro, algo que no puede hacer un planeta oceánico con nubes y
algo de hielo. La otra salida es un efecto invernadero más fuerte.
\end{solution}

\begin{solution}{exo:g12:early-earth:13}
(1) $273 - 237 = \qty{36}{K}$ para llegar a la congelación; $288 - 237 = \qty{51}{K}$ para llegar a la media actual. (2) El efecto invernadero actual
añade $288 - 255 =
\qty{33}{K}$: la Tierra joven necesitaba uno más fuerte.
(3) \qty{0.1}{bar} es una décima parte de la presión de toda la atmósfera
actual, solo para el dióxido de carbono.
\end{solution}

\begin{solution}{exo:g12:early-earth:14}
(1) Antes de hace 2450 millones de años la firma existe, así que el aire
contenía menos de $10^{-5}$ del oxígeno actual; después de hace 2090 millones
de años falta, así que el aire contenía más. (2) Sí: el \omterm{def:g12:early-earth:goe}{Gran Evento de
Oxidación} va de hace unos 2460 a 2060 millones de años, el mismo intervalo.
\end{solution}

\begin{solution}{exo:g12:early-earth:15}
(1) Hace unos $4567 - 60 = 4507$ millones de años. (2) Como máximo $4507 - 4404 =
103$ millones de años después del impacto: la Tierra tenía una corteza y agua
en el plazo de unos cien millones de años tras el impacto.
\end{solution}

\begin{solution}{pb:g12:early-earth:1}
\textbf{1.} $4.567 - 4.0 = 0.567$ miles de millones de años.

\textbf{2.} $1 - 0.567/4.567 \approx 1 - 0.124 = 0.876$.

\textbf{3.} $L = 1/(1 + 0.4 \times 0.876) = 1/1.350 \approx 0.741$ de la
actual.

\textbf{4.} $0.741 \times 1361 \approx \qty{1008}{W/m^2}$.

\textbf{5.} $1 - 0.741 \approx 0.26$: un \qty{26}{\percent} menos brillante.

\textbf{6.} $T_{\mathrm{eq}} = \left(S(1-A)/(4\sigma)\right)^{1/4}$.

\textbf{7.} Con $A$ y $\sigma$ fijos, $T_{\mathrm{eq}}$ es una constante por
$S^{1/4}$, así que $T'/T = (S'/S)^{1/4}$.

\textbf{8.} $0.741^{1/4} = \sqrt{\sqrt{0.741}} \approx \sqrt{0.861} \approx
0.928$.

\textbf{9.} $255 \times 0.928 \approx \qty{237}{K}$.

\textbf{10.} $237 - 273 \approx \qty{-36}{\celsius}$.

\textbf{11.} Que había agua líquida en la superficie, un océano, hace 4400
millones de años.

\textbf{12.} En aguas poco profundas, como los \omterm{def:g12:early-earth:stromatolite}{estromatolitos} vivos de hoy.

\textbf{13.} Las \omterm{def:g12:early-earth:bif}{formaciones de hierro bandeado} de los montes Hamersley,
depositadas en un fondo marino.

\textbf{14.} Sí: \qty{237}{K} está muy por debajo de \qty{273}{K}; con la
atmósfera actual, la superficie media habría estado congelada.

\textbf{15.} Las rocas muestran agua líquida bajo un Sol demasiado débil para
mantenerla líquida con la atmósfera actual: la \omterm{def:g12:early-earth:faint-young-sun}{paradoja del Sol joven débil}.

\textbf{16.} $288 - 255 = \qty{33}{K}$.

\textbf{17.} Al menos $273 - 237 = \qty{36}{K}$, más que el actual.

\textbf{18.} $A' = 1 - 0.706/0.741 \approx 0.047$: un planeta casi negro, no
es verosímil.

\textbf{19.} El dióxido de carbono, mucho más abundante que hoy; el metano,
producido por los primeros microbios, pudo ayudar.

\textbf{20.} La \omterm{def:g10:habitable-planet:equilibrium-temperature}{temperatura de equilibrio} de la Tierra hace cuatro mil
millones de años era de \textbf{unos 237 K}, unos \qty{36}{\celsius} por
debajo del punto de congelación.
\end{solution}
